This companion result extends the domain of an existing positive-variable
construction. Its numerical saving is only two operations; its new
argument proves that a formerly excluded congruence class cannot create
a spurious solution.

This section uses the separate, coherent compact44 tuple
$\beta_+=2^{590069}-1$, $K_0,M_C,M_F$, with $M_F$ already shifted.
It is not interchangeable with Section~\ref{sec:integer}'s source-34
tuple. The four exact payloads and their finite recipe are retained under
\texttt{positive/coefficients/}.

\subsection{The complete positive-variable polynomials}
For twelve strictly positive coordinates use
\[
 J,F,\alpha,t,T,X,Y,r,\eta,\zeta,n,\gamma.
\]
For thirteen use
\[
 J,F,\alpha,f,i,T,X,Y,r,\eta,\zeta,u,\gamma.
\]
All symbols in this subsection are local to these formulas. Define
\begin{align*}
 q&=\beta_+J+1,&C&=q-F-\alpha,&E&=XY,&k&=\eta+\zeta,\\
 c&=kY+\eta,&a&=(X+1)Y,&H&=4a+3,&\Delta&=(a+2)^2-1,\\
 \ell&=EkY,&M&=q^3(q-1),&W&=(q-1)k+1,&U&=q^2Y,\\
 B_t&=(qK_0+X)C-qF,&&&&&
\end{align*}
and
\begin{align*}
 Z&=M(\ell+k+r)+(q-1)Y-q^3WB_t,\\
 A_0&=Z^2-M^2[\ell(\ell+k)+1],\\
 R&=[(q-1)(q-F)+\alpha](q^2-1)+(M_C+qM_F)J,\\
 v&=k-R-1,\\
 D_*&=U(X+ac)+H(U\gamma+qE-v),\\
 A_1&=D_*^2-U^2(1+\Delta c^2).
\end{align*}
In both tails put $Q=\Delta(f^2-1)$ and $V=f^2(cT-R)-c$.
In the twelve-coordinate case $f=t+1$ is an abbreviation. Let
\[
 d=c^4,\quad h_0(z)=(1-Q)z^2+2(1-Q)Vz+V^2-1,
\]
\[
 b_n=2(1-Q)(n+V),\quad A_n=d\,h_0(n)+(1-Q)Q,\quad
 L=A_n^2-d\,b_n^2Q.
\]
The twelve-positive equation is
\[
 P_{12}^+=(1+A_0)(1+A_1)(1+L)-1=0.
\]
For thirteen positive coordinates set $I=ic^2$ and $y_*=V+u$, giving
\[
 P_{13}^+=(1+A_0)(1+A_1)
       [QV^2+(1-Q)y_*^2][1+I^2-Q]-1=0.
\]
Their complete DAGs have 98 and 88 operations and exact degrees 70 and 51.
The parents first compute $J=6j+1$ or $J=3j+1$; the new versions accept
positive $J$ directly and remove those two operations.

\subsection{Recovering divisibility despite a common factor three}
The product finalizers still recover their intended equations without
assuming $q$ even. Indeed $M=q^3(q-1)$ is always even;
$\Delta\equiv0$ or $3\pmod4$, so $A_1$ is never $2\pmod4$.
The remaining tail factors have the same exclusion of $-1$ as in the
ordinary parent proofs. Thus their unit argument applies on the enlarged
positive domain.

The unchanged first carrier proves $C>0$, $q^3\mid Y$, $q\ge16$,
and $3q<R<q^4-2$. The main carrier then gives a positive integer
$D=D_*/U$ with $D^2-\Delta c^2=1$ and
\[
 U\mid H(qE-v).
\]
The needed weaker coprimality statement is exact:
\begin{equation}
 \gcd(q,H)=\gcd(q^5,H)=\gcd(q,3)\in\{1,3\}.\label{eq:gcd3}
\end{equation}
For primes other than three this follows from $H\equiv3\pmod{q^3}$.
If $3\mid q$, then $27\mid Y$ and $H\equiv3\pmod{27}$, so its
$3$-adic valuation is exactly one. Since $q^5\mid U$, the displayed
divisibility still yields $q^4\mid(qE-v)$, and $q^4\mid qE$ gives
$q^4\mid v$.

The retained first/main Pell classification supplies indices $n_0,p$ with
\[
 k=2\psi_{2XY^2+1}(n_0),\quad c=\psi_{a+2}(p),\quad p>n_0,
 \quad 2n_0\equiv R+1\pmod{q^4}.
\]
Consequently $n_0\ge25$. The early rank bounds and the parent ordinary
strong and auxiliary arguments, now licensed by $q^4\mid v$, give
$p=R$ and then $R=2n_0-1$ exactly. These arguments do not use the
discarded residue restriction. In particular $U\mid v$, and the main
carrier gives
\begin{equation}
 q\mid3X,\qquad X\equiv2^R\pmod{H/g},\quad
 g=\gcd(q,H)\in\{1,3\}.\label{eq:weakprojection}
\end{equation}
The new task is to turn this weaker projection into equality before
assuming that $q$ divides $X$.

\subsection{Pell growth eliminates the exceptional branch}
Write $R=2r_0+1$, with $r_0=n_0-1\ge24$, and $P_0=2XY^2+1$,
$A=(X+1)Y+2$. Pell growth gives
\[
 \frac ck\ge\frac12\left(\frac{(2A-1)^2}{2P_0}\right)^{r_0},
 \qquad\frac ck<Y+1.
\]
The exact identity
\[
 X(2A-1)^2-2P_0(X+1)^2
 =12XY(X+1)+9X-2(X+1)^2>0
\]
for $X,Y\ge1$ improves this to
\begin{equation}
 Y+1>\frac12\left(\frac{(X+1)^2}{X}\right)^{r_0}.\label{eq:growth}
\end{equation}
For $X\ge2$, the base is at least $9/2$. The elementary inequality
$9^{16}>6\cdot8^{16}$ and $r_0\ge24$ imply $Y>2^R$.
Both $X$ and $2^R$ lie in $(0,a)$, whereas $H/g\ge H/3>a$.
Their congruence in \eqref{eq:weakprojection} therefore gives $X=2^R$.

For $X=1$, equation~\eqref{eq:growth} implies $Y\ge2^{R-2}$.
Then $H=8Y+3$ and $0<2^R-1<H/2$. Congruence modulo $H$ is impossible;
modulo $H/3$, the only possible positive multiple is $H/3$ itself.
It would give $8Y=3\cdot2^R-6$, impossible modulo eight for $R\ge3$.
Thus $X=1$ is excluded and again $X=2^R$.

Now $q\mid3\cdot2^R$ makes $q=2^u$ or $q=3\cdot2^u$.
The authentic Mersenne modulus eliminates the latter. If
$2^d-1\mid3\cdot2^u-1$, reduce $u$ modulo $d$ to $r\in\{0,\ldots,d-1\}$.
For $r\le d-2$, the number $3\cdot2^r$ lies strictly between one and
$2^d-1$. For $r=d-1$, its residue is $2^{d-1}+1$, also not one.
Here $d=590069\ge4$. Hence $q$ is a power of two and $q\mid X$.
The original source witnesses can now be restored exactly as in the
parent. Canonical refutation witnesses already have positive $J$, so
completeness is unchanged. This proves the same ordinary consistency
equivalence on the larger domain.

The exact degree certificates remain nonzero after replacing the old
leading coefficient $6\beta_+$ or $3\beta_+$ of $q$ by $\beta_+$.
The retained leading coefficients are $64\beta_+^{18}$ and
$-32\beta_+^{15}$. The circuit check confirms the counts
$54M+25A+19S=98$ and $48M+23A+17S=88$.

